\documentclass[11pt,reqno]{amsart}

\usepackage[T1]{fontenc}
\usepackage[utf8]{inputenc}
\usepackage{lmodern}
\usepackage{microtype}
\usepackage{amsmath,amssymb,mathtools}
\usepackage{booktabs}
\usepackage{enumitem}
\usepackage{xcolor}
\usepackage{hyperref}
\usepackage[nameinlink,capitalise,noabbrev]{cleveref}

\hypersetup{
  colorlinks=true,
  linkcolor=blue!55!black,
  citecolor=blue!55!black,
  urlcolor=blue!55!black,
  pdftitle={A log-free lower bound for imaginary quadratic fields of 5-rank at least two},
  pdfauthor={Hao Zhang}
}

\newtheorem{theorem}{Theorem}[section]
\newtheorem{proposition}[theorem]{Proposition}
\newtheorem{lemma}[theorem]{Lemma}
\newtheorem{corollary}[theorem]{Corollary}
\theoremstyle{definition}
\newtheorem{definition}[theorem]{Definition}
\theoremstyle{remark}
\newtheorem{remark}[theorem]{Remark}

\newcommand{\Q}{\mathbb{Q}}
\newcommand{\Z}{\mathbb{Z}}
\newcommand{\F}{\mathbb{F}}
\newcommand{\Cl}{\operatorname{Cl}}
\newcommand{\Disc}{\operatorname{Disc}}
\newcommand{\Jac}{\operatorname{Jac}}
\newcommand{\rk}{\operatorname{rk}}
\newcommand{\Ht}{\operatorname{H}}
\newcommand{\PGL}{\operatorname{PGL}}
\newcommand{\SL}{\operatorname{SL}}
\newcommand{\NoLogCount}{N^{-}_{5,2}}

\title[A log-free lower bound for quadratic fields]{A log-free lower bound for imaginary quadratic fields of 5-rank at least two}
\author{Hao Zhang}
\date{Research manuscript draft, 14 August 2026}

\begin{document}

\begin{abstract}
Let \(\NoLogCount(X)\) denote the number of imaginary quadratic fields
\(K\), counted up to \(\Q\)-isomorphism, for which
\(\lvert\Disc(K)\rvert\leq X\) and
\(\dim_{\F_5}\Cl(K)/5\Cl(K)\geq 2\).  Bartz, Levin, and Thamminana
proved the lower bound
\(\NoLogCount(X)\gg X^{1/3}/(\log X)^2\).  We present a complete
human-readable proof candidate for the log-free estimate
\[
  \NoLogCount(X)\gg X^{1/3}.
\]
The argument combines their explicit genus-two curve with the strong
power-free-value theorem of Stewart--Top, rather than the logarithmically
weaker variant used in the Hilbert-irreducibility framework of
Kulkarni--Levin.  The main points are to preserve one positive congruence
cone, extract bounded witnesses for distinct squarefree values, delete a
thin set without losing the quadratic main term, and identify the displayed
genus-two curve with the model to which the torsion theorem applies.  We give
the required endpoint-avoidance and square-class injection lemmas, an exact
rational isomorphism of the two genus-two models, and a reproducibility
statement.  The full analytic argument has undergone independent adversarial
review, but has not been formalized end to end in a proof assistant; this
document should therefore be treated as a pre-submission research draft.
\end{abstract}

\maketitle

\section{Introduction}

For a positive integer \(m\) and a finite abelian group \(A\), write
\[
  \rk_m A=\max\{r:(\Z/m\Z)^r\hookrightarrow A\}.
\]
When \(m=p\) is prime, this is
\(\dim_{\F_p}(A/pA)\).  We study the counting function
\begin{equation}\label{eq:counting-function}
  \NoLogCount(X)
  =\#\left\{
  \begin{array}{c}
  K/\Q\text{ imaginary quadratic, up to }\Q\text{-isomorphism}:\\
  |\Disc(K)|\leq X,\quad \rk_5\Cl(K)\geq 2
  \end{array}
  \right\}.
\end{equation}

Byeon proved the unconditional estimate
\(\NoLogCount(X)\gg X^{1/4}\) \cite{Byeon2009}; an earlier result gives a
weaker exponent in greater generality \cite{Byeon2006}.  Kulkarni and Levin
developed a Hilbert-irreducibility method that converts rational torsion on
hyperelliptic Jacobians into quantitative class-group constructions
\cite{KulkarniLevin2022}.  Bartz, Levin, and Thamminana applied this framework
to an explicit genus-two curve and obtained
\begin{equation}\label{eq:blt-bound}
  \NoLogCount(X)\gg \frac{X^{1/3}}{(\log X)^2}
  \qquad\text{\cite[Theorem~1.3]{BartzLevinThamminana2025}}.
\end{equation}
Conditional results of a different shape are known; for example,
Chattopadhyay and Saikia obtain, assuming the \(abc\)-conjecture, an exponent
\(1/4-\varepsilon\) when \(p=5\) \cite{ChattopadhyaySaikia2023}.

The purpose of this paper is to remove the logarithmic loss in
\eqref{eq:blt-bound} without changing the genus-two input.

\begin{theorem}\label{thm:main}
There exist constants \(c>0\) and \(X_0>1\) such that, for every
\(X\geq X_0\),
\[
  \NoLogCount(X)\geq cX^{1/3}.
\]
\end{theorem}

The logarithmic factor in \eqref{eq:blt-bound} comes from the power-free-value
estimate used in the general Kulkarni--Levin construction.  For the particular
binary sextic arising here, the stronger theorem of Stewart and Top
\cite[Theorem~1]{StewartTop1995} applies.  Its proof produces
\(\gg H^2\) distinct squarefree values from a fixed positive box of side
\(H\), whereas the relevant thin set contains only \(O(H)\) rational
parameters of comparable height.  The difference remains quadratic.

There are three bookkeeping issues that must be resolved before this idea
becomes a proof.  First, a projective change of variable in the
Stewart--Top proof may change the positive cone; we avoid it by choosing the
Kulkarni--Levin coordinate so that both endpoint coefficients are nonzero.
Second, the thin-set estimate counts rational parameters rather than binary
form values; even homogeneity gives the required injection.  Third, the
genus-two model in the torsion construction must be identified with the
displayed curve by an exact rational map.  These points are treated in
\Cref{sec:curve,sec:localized,sec:main-proof}.

To the best of our knowledge within a search of OpenAlex, Semantic Scholar,
Crossref, arXiv, and the direct citation chain completed on 14 August 2026,
no earlier source states the unconditional log-free estimate in
\Cref{thm:main}.  The closest unconditional result with the same exponent is
\eqref{eq:blt-bound}.  This is a scoped literature statement, not a claim that
all indexed and unindexed literature has been exhausted.

\section{The fixed genus-two curve}\label{sec:curve}

Bartz, Levin, and Thamminana exhibit the curve
\begin{equation}\label{eq:C0}
  C_0:\quad y^2=f_0(x)
  =(5x+7)(128x^4+549x^3+1007x^2+936x+368)
\end{equation}
and obtain the needed rational \(5\)-torsion input from their
Theorem~2.1 \cite[pp.~4--7]{BartzLevinThamminana2025}.  The following form of
that input is all that we use.

\begin{proposition}\label{prop:C0-rank}
The curve \(C_0\) has genus two, and
\[
  \rk_5\Jac(C_0)(\Q)_{\mathrm{tors}}\geq 2.
\]
Moreover, after the change of variables \(u=640x\), \(v=640^2y\), it has the
monic odd-degree model
\begin{equation}\label{eq:g-model}
\begin{aligned}
  v^2=g(u)=(u+896)\bigl(&u^4+2745u^3+3222400u^2\\
  &+1916928000u+482344960000\bigr).
\end{aligned}
\end{equation}
The associated projective branch form has irreducible-factor degrees
\(1,1,4\) over \(\Q\).
\end{proposition}

\begin{proof}
The identity \(g(u)=640^4f_0(u/640)\) is immediate by expansion.  The quartic
factor in \eqref{eq:g-model} is irreducible over \(\Q\), and \(g\) is
separable.  Together with the branch point at infinity, this gives factor
degrees \(1,1,4\).

For the torsion assertion, substitute
\((t,u,z)=(2/3,-1/3,25)\) into the coefficient formulas of
\cite[Theorem~2.1]{BartzLevinThamminana2025}.  The resulting even model
\(C_{\mathrm{even}}:y^2=P(x)\) is mapped to \(C_0\) by
\begin{equation}\label{eq:forward-map}
  x_0=\frac{2(1-2x)}{5x-1},\qquad
  y_0=\frac{59049y}{1024(5x-1)^3},
\end{equation}
with inverse
\begin{equation}\label{eq:inverse-map}
  x=\frac{x_0+2}{5x_0+4},\qquad
  y=\frac{8192y_0}{2187(5x_0+4)^3}.
\end{equation}
The exact polynomial identity and the two coordinate compositions are recorded
in \Cref{app:model}.  Thus the theorem-level torsion conclusion for the even
model applies to \(C_0\).  The stronger reported computation of the entire
torsion subgroup is unnecessary here.
\end{proof}

\section{Two elementary bridge lemmas}\label{sec:bridge}

The first lemma lets us preserve the positive cone in the later
Stewart--Top count.

\begin{lemma}[Finite endpoint avoidance]\label{lem:endpoint}
Fix the bad-prime data \(M_0\), a shift \(N_0/M_0\), and the polynomial
\(g\) in \eqref{eq:g-model}.  In the proof of
\cite[Lemma~3.1]{KulkarniLevin2022}, the integer \(N_1\) may be chosen
sufficiently large and so that
\begin{equation}\label{eq:endpoint-avoid}
  g\!\left(-\frac{N_0}{M_0}+\frac1{N_1}\right)
  g\!\left(-\frac{N_0}{M_0}-\frac1{N_1}\right)\neq 0.
\end{equation}
For this choice, the transformed projective branch form has nonzero leading
and trailing coefficients.
\end{lemma}

\begin{proof}
The construction of Kulkarni--Levin permits any sufficiently large positive
integer \(N_1\).  Each equality obtained by setting one factor in
\eqref{eq:endpoint-avoid} equal to zero excludes only finitely many such
integers, because \(g\) has finitely many roots.  The two values in
\eqref{eq:endpoint-avoid} are precisely the images of the projective endpoints
under the inverse M\"obius transformation used below.  Their nonvanishing is
equivalent to the two endpoint coefficients being nonzero.
\end{proof}

\begin{lemma}[Even-form square-class injection]\label{lem:injection}
Let \(F\in\Z[X,Y]\) be homogeneous of even degree six, and fix
\(w\in\Z\setminus\{0\}\).  Suppose
\[
  F(a_i,b_i)=t_iw^2\qquad(i=1,2),
\]
where \(t_1,t_2\) are signed squarefree integers and \(b_1b_2\neq0\).  If
\(a_1/b_1=a_2/b_2\), then \(t_1=t_2\).  The same conclusion holds after a
fixed invertible fractional-linear transformation of the ratios.
\end{lemma}

\begin{proof}
Equality of the ratios gives
\((a_2,b_2)=\lambda(a_1,b_1)\) for some \(\lambda\in\Q^\times\).  Hence
\[
  \frac{t_2}{t_1}
  =\frac{F(a_2,b_2)}{F(a_1,b_1)}
  =\lambda^6=(\lambda^3)^2.
\]
Two signed squarefree integers representing the same element of
\(\Q^\times/\Q^{\times2}\) are equal.  An invertible fractional-linear
transformation preserves equality of ratios.
\end{proof}

\section{A localized log-free squarefree-value proposition}\label{sec:localized}

We isolate the exact consequence of Stewart--Top needed in the proof.  The
statement includes the bounded positive witnesses that are present in the
proof of their theorem, although not displayed in the definition of their
unbounded counting function.

\begin{proposition}\label{prop:localized-ST}
Let \(F\in\Z[X,Y]\) be a squarefree binary form of degree \(r=6\), with
nonzero leading and trailing coefficients.  Suppose the largest degree of an
irreducible factor over \(\Q\) is \(m\leq 5\).  Fix integers \(A,B,M\), with
\(M\geq1\), and let \(w\geq1\) be maximal such that
\[
  w^2\mid F(a,b)
  \quad\text{whenever}\quad
  a\equiv A\pmod M,\quad b\equiv B\pmod M.
\]
Then there are constants \(c_0>0\) and \(H_0\geq1\), depending only on the
fixed data, such that every real \(H\geq H_0\) admits at least \(c_0H^2\)
distinct signed squarefree integers \(t\) for which
\[
  F(a,b)=tw^2
\]
has a solution with
\(1\leq a,b\leq H\),
\(a\equiv A\pmod M\), and
\(b\equiv B\pmod M\).
\end{proposition}

\begin{proof}
Apply \cite[Theorem~1]{StewartTop1995} with \(k=2\).  Its hypothesis on the
factor degrees is \(m\leq2k+1=5\).  The bounded-witness conclusion follows
from the proof: the set constructed in
\cite[pp.~951--952]{StewartTop1995} lies in the positive box
\(1\leq a,b\leq H\) and in the fixed residue class.  Equation~(10) on
p.~953 gives \(\gg H^2\) eligible pairs.  Equations~(11)--(15), together with
the uniform Thue-equation estimate of \cite[Corollary~1]{Stewart1991}, bound
the number of primitive representations of any one value by a constant
depending only on the fixed form and residue data.  The paragraph following
(15) therefore produces \(\gg H^2\) distinct values, not merely pairs.

The cited proof is first written for sufficiently large integral box
parameters.  For real \(H\), use \(\lfloor H\rfloor\geq H/2\) after enlarging
\(H_0\), and shrink the constant by a factor of four.
\end{proof}

\begin{remark}\label{rem:source-bookkeeping}
The proof of Proposition~\ref{prop:localized-ST} uses the distinct-value set obtained
after the representation-multiplicity estimate, rather than identifying it
with the earlier pair set.  Also, a preliminary \(\SL_2(\Z)\) change of
variables in the general proof is unnecessary once the two endpoint
coefficients are nonzero.  This is why Lemma~\ref{lem:endpoint} is imposed before
the final congruence class is fixed.
\end{remark}

\section{The class-group specialization input}\label{sec:KL}

We use the following specialized form of the Kulkarni--Levin construction.
It is a direct combination of their local-approximation lemma, thin-set height
bound, specialization theorem, and the proof of their imaginary quadratic
case \cite[Lemmas~3.1--3.2, Theorems~1.3 and~4.1]{KulkarniLevin2022}.

Fix a sufficiently large positive integer \(N_0\), coprime to the product
\(M_0\) of the bad primes of \(C_0\), such that
\begin{equation}\label{eq:negative-neighborhood}
  g(u)<0\qquad\text{if}\qquad
  \left|u+\frac{N_0}{M_0}\right|<1.
\end{equation}
This is possible because \(g\) is monic of odd degree.  Put
\(\phi(P)=u(P)+N_0/M_0\).  The specialization theorem supplies a fixed thin
set \(\Omega\subset\Q\) outside which the relevant torsion classes specialize
to ideal classes.

Choose \(N_1\) as in Lemma~\ref{lem:endpoint} and define
\begin{equation}\label{eq:mobius}
  \psi(\alpha)=\frac{1-N_1\alpha}{1+N_1\alpha},
  \qquad
  \tau(q)=\psi^{-1}(q)=\frac{1-q}{N_1(1+q)}.
\end{equation}
The proof of \cite[Lemma~3.1]{KulkarniLevin2022} then gives fixed integers
\(A,B,M\)---one may take \(A=B=1\)---such that positive integers in this
single residue class satisfy all prescribed local conditions.  The
composition order is
\begin{equation}\label{eq:composition-order}
  q=(\psi\circ\phi)(P)=a/b,
  \qquad
  \phi(P)=\tau(a/b),
  \qquad
  u(P)=\tau(a/b)-N_0/M_0.
\end{equation}

After clearing fixed denominators and removing fixed square factors, the
branch equation can be written
\begin{equation}\label{eq:binary-form-identity}
  g\!\left(\tau(X/Y)-\frac{N_0}{M_0}\right)
  =F(X,Y)R(X,Y)^2,
\end{equation}
where \(F\in\Z[X,Y]\) is a squarefree binary sextic and
\(R\in\Q(X,Y)\).  The projective change of coordinate preserves separability
and factor degrees, so \(F\) has factor-degree pattern \(1,1,4\).  By
Lemma~\ref{lem:endpoint}, its leading and trailing coefficients are nonzero.

For the final residue class define \(w\) to be the maximal positive integer
whose square divides every value in that class, exactly as in
\cite[p.~948]{StewartTop1995}.  All of
\(N_0,M_0,N_1,A,B,M,F,w,\Omega\) are now fixed; none depends on the later
height parameter.

\section{Proof of the main theorem}\label{sec:main-proof}

\begin{proof}[Proof of \Cref{thm:main}]
Apply Proposition~\ref{prop:localized-ST} to the binary sextic \(F\) from
\eqref{eq:binary-form-identity}.  There are \(\gg H^2\) distinct signed
squarefree integers \(t\), each with a chosen witness \((a_t,b_t)\) in the
same positive box and residue class, such that
\begin{equation}\label{eq:Ftw}
  F(a_t,b_t)=t w^2.
\end{equation}
The constants are uniform for every sufficiently large \(H\).

For each \(t\), set \(q_t=\tau(a_t/b_t)\).  A fixed fractional-linear map
distorts multiplicative height by only a fixed constant, hence
\begin{equation}\label{eq:height-distortion}
  \Ht(q_t)\ll H.
\end{equation}
By Lemma~\ref{lem:injection}, the map \(t\mapsto q_t\) is injective.  The thin-set
bound \cite[Lemma~3.2]{KulkarniLevin2022} therefore shows that only \(O(H)\)
of the chosen parameters belong to \(\Omega\).  Removing them leaves
\begin{equation}\label{eq:after-deletion}
  \gg H^2-O(H)\gg H^2
\end{equation}
distinct squarefree values.

For a surviving value, the fixed real neighborhood
\eqref{eq:negative-neighborhood} and
\eqref{eq:binary-form-identity} imply \(t<0\), since the omitted factors and
\(w^2\) are squares.  The field computation in the proof of
\cite[Theorem~1.2, equations~(3.1)--(3.2)]{KulkarniLevin2022} gives
\begin{equation}\label{eq:field-identification}
  \Q(P_t)=\Q(\sqrt{F(a_t,b_t)})=\Q(\sqrt t).
\end{equation}
The local choices and the specialization theorem, together with
Proposition~\ref{prop:C0-rank}, yield
\begin{equation}\label{eq:rank-conclusion}
  \rk_5\Cl(\Q(\sqrt t))
  \geq \rk_5\Jac(C_0)(\Q)_{\mathrm{tors}}
  \geq2.
\end{equation}
Distinct signed squarefree \(t\)'s give distinct quadratic fields.  In
particular, the two conjugate points over one rational parameter are not
double-counted: the argument counts squarefree values \(t\), not geometric
points.

Since \(F\) is fixed and has degree six, there is a constant \(C_F>0\) such
that
\[
  |F(a,b)|\leq C_FH^6
  \qquad(1\leq a,b\leq H).
\]
The fundamental discriminant of \(\Q(\sqrt t)\) has absolute value at most
\(4|t|\).  From \eqref{eq:Ftw} we obtain a fixed \(C_\Delta>0\) with
\begin{equation}\label{eq:disc-height}
  |\Disc(\Q(\sqrt t))|\leq C_\Delta H^6.
\end{equation}
Given sufficiently large \(X\), take
\[
  H=\left\lfloor (X/C_\Delta)^{1/6}\right\rfloor.
\]
Then \(H\gg X^{1/6}\), all the surviving fields satisfy the discriminant
cutoff \(X\), and their number is \(\gg H^2\gg X^{1/3}\).  This proves the
stated estimate.
\end{proof}

\section{Quantifiers, dependencies, and scope}\label{sec:scope}

The order of choices in the proof is
\begin{equation}\label{eq:quantifier-order}
\begin{aligned}
  C_0
  &\Longrightarrow(S,M_0,N_0,\phi,\Omega)
  \Longrightarrow N_1\\
  &\Longrightarrow(A,B,M,F,w,c_0,H_0,C_\Delta)\\
  &\Longrightarrow \forall H\geq H_0
  \Longrightarrow \forall X\geq X_0.
\end{aligned}
\end{equation}
Thus the residue class, the class-wise fixed square divisor, the thin set, all
exceptional sets, and every implied constant are selected before \(H\) and
\(X\).  This prevents a variable local condition or a height-dependent square
factor from being hidden in the Vinogradov symbol.

Three source-level bookkeeping points are worth recording.
\begin{enumerate}[label=(\roman*)]
  \item If a binary form is precomposed with a matrix \(L\), the residue class
  preserving the original values is transported by \(L^{-1}\).  In the
  present argument the optional change is not used, because
  Lemma~\ref{lem:endpoint} makes both endpoint coefficients nonzero.
  \item The pair set in the proof of \cite[Theorem~1]{StewartTop1995} is not
  itself the distinct-value set.  The passage to distinct values occurs only
  after the representation-multiplicity estimates (11)--(15).
  \item Two conjugate points can generate the same quadratic field.  Counting
  distinct signed squarefree values \(t\) removes this multiplicity.
\end{enumerate}
These observations do not change the exponent, but they are needed to keep the
objects being counted aligned throughout the proof.

\section{Related work and novelty boundary}\label{sec:related}

The result closest to \Cref{thm:main} is the unconditional estimate
\eqref{eq:blt-bound}.  Byeon's \(X^{1/4}\) theorem is log-free but has a
smaller exponent \cite{Byeon2009}.  The general Kulkarni--Levin corollary for
\(m=5\) has exponent \(1/5\) and a \((\log X)^{-2}\) factor
\cite[Corollary~1.4]{KulkarniLevin2022}.  The conditional result of
Chattopadhyay--Saikia has exponent \(1/4-\varepsilon\) for \(p=5\)
\cite{ChattopadhyaySaikia2023}.

The scoped search supporting this comparison included the exact notations
\(\NoLogCount(X)\) and \(N^-(5^2;X)\), equivalent finite-abelian-group
formulations, citation descendants of the principal sources, and current
records in OpenAlex, Semantic Scholar, Crossref, and arXiv through 14 August
2026.  It did not exhaust subscription databases, nonpublic manuscripts,
unindexed theses, non-English literature, or work affected by indexing lag.
Accordingly, the appropriate scholarly formulation is ``to the best of our
knowledge within the stated search scope,'' rather than an absolute priority
claim.

\section{Verification and reproducibility}\label{sec:reproducibility}

The proof was assembled from frozen versions of the four principal sources
\cite{BartzLevinThamminana2025,KulkarniLevin2022,StewartTop1995,Stewart1991}
and subjected to independent adversarial review.  Exact rational arithmetic
checks the model identity in Proposition~\ref{prop:C0-rank}, the forward curve identity,
both coordinate compositions, and the rejection of the incorrect inverse
coefficient \(1024\), which would return \(y/8\) instead of \(y\).

In the accompanying project, the principal replay commands are
\begin{verbatim}
python3 mathematics/worker/no-log-blt-c0-rank-replay.py
python3 mathematics/worker/verify-route-audit.py
lean -q -t 0 mathematics/formal/NoLogSourceBridge.lean
lean -q -t 0 mathematics/formal/NoLogInference.lean
\end{verbatim}
The two Lean files certify only the final conditional numerical assembly and
elementary algebraic components.  They do not formalize the Stewart--Top,
Kulkarni--Levin, or Bartz--Levin--Thamminana inputs, nor the full unconditional
analytic bridge.  The complete theorem is therefore not claimed here as an
end-to-end machine-verified result.  The manuscript source and checks are bound
to project commit
\texttt{c9ca34210ee58d01c1f6fc5cd97f969534dbaa70} before the manuscript
commit itself.

\section*{AI assistance disclosure}

DeepScientist, an AI-assisted research system, was used for structured
literature discovery, extraction of proof obligations, symbolic-algebra
replay, adversarial consistency checking, and preparation of the initial
LaTeX draft.  DeepScientist is not an author.  Hao Zhang is the sole author,
retains responsibility for every mathematical statement and citation, and
should complete the conventional final human source and proof review required
by the intended journal before submission.

\section*{Acknowledgements}

The author thanks the authors of the cited works for making the constructions
and source materials on which this argument builds available.

\appendix

\section{Exact identification of the two genus-two models}\label{app:model}

For the sample \((t,u,z)=(2/3,-1/3,25)\) in
\cite[Section~2]{BartzLevinThamminana2025}, exact substitution gives
\[
\begin{aligned}
  a&=\frac{2048}{243},
  &a_0&=-\frac{10240}{531441},
  &a_2&=\frac{20480}{59049},\\
  a_4&=\frac{182272}{177147},
  &a_6&=\frac{31129600}{531441}.
\end{aligned}
\]
Let
\[
  P(x)=a(a_6x^6+a_4x^4+a_2x^2+a_0).
\]
Writing \(D=5x-1\) and \(N=2(1-2x)\), the forward map
\eqref{eq:forward-map} is certified by the coefficientwise identity
\begin{equation}\label{eq:forward-identity}
  1024^2D^6 f_0(N/D)=59049^2P(x).
\end{equation}
Moreover,
\[
  5x_0+4=\frac6D,
  \qquad
  \frac{x_0+2}{5x_0+4}=x,
\]
and the \(y\)-coordinate composition factor is
\[
  \frac{8192}{2187(5x_0+4)^3}
  \frac{59049}{1024D^3}=1.
\]
The old numerator \(1024\) in the first factor would give \(1/8\), so the
coefficient \(8192\) in \eqref{eq:inverse-map} is forced.  These are rational
function identities on the dense affine opens where the denominators are
nonzero; the induced maps extend to the smooth projective curves.

\bibliographystyle{amsplain}
\bibliography{log-free-5-rank}

\end{document}
